\documentclass{article}
\usepackage{/home/user1/code/tex/sty}
\usepackage{amsfonts}
\usepackage{amsmath}
\usepackage{geometry}
\usepackage{graphicx}
\usepackage{hyperref}
\setlength{\parindent}{0pt}
\geometry{top=1in,bottom=1in,left=1in,right=1in}
\usepackage{amstext}
\usepackage{array}
\newcolumntype{L}{>{$}l<{$}}
\newcommand{\ctab}[2]{\begin{center}\begin{tabular}{#1}#2\end{tabular}\end{center}}
\newcommand{\ctabt}[2]{\begin{center}\begin{tabular}{#1}#2\end{tabular}\end{center}}
\n{rb}{\ r_{\textrm{Bohr}}}
\n{taup}{\tau_{\textrm{proton}}}
\n{tmag}{\ T_\textrm{magnetar}}
\n{tro}{T_\textrm{room}}
\n{tcmb}{T_\textrm{CMB}}
\n{tsun}{T_\textrm{sun}}
\n{tuni}{t_\textrm{Universe}}
\n{rhouni}{\rho_\textrm{Universe}}
\n{rhomw}{\rho_\textrm{Milky}}
\n{rhoair}{\rho_\textrm{air}}
\n{rhowater}{\rho_\textrm{water}}
\n{rhoau}{\rho_\textrm{Au}}
\n{me}{\textrm{ m}}
\n{se}{\textrm{ sec}}
\n{yes}{\rom{Yes}}
\n{no}{\rom{No}}
\n{ev}{\textrm{ eV}}
\n{psb}{\bar{\psi}}
\n{sib}{\bar{\sigma}}
\n{ua}{U(1)_A}
\n{uv}{U(1)_V}
\begin{document}
\begin{center}
{\Large{Problem Set 1}}\\~\\
\begin{tabular}{l l}
Student&Agustin Romero\\
Instructor&Savas Dimopolous\\
Assistant&Jed Oliver\\
Course&Standard Model\\    
Institution&Stanford University\\
Date&\today\\
\end{tabular}
\end{center}
\section*{Problem 1.a.}
These are the variables converted using only $\hbar=1.0545\ten{-34}\un{ Js}$ and $c=299792458\textrm{ m}\textrm{ s}^{-1}$.
\ctab{L L L L}{\text{Quantity} & \text{eV} & \text{m} & \text{s}\\
1\un{GeV} & 1\ten{9}\ev & 5\ten{15}\me^{-1}&2\ten{24}\se^{-1}\\
1\un{sec}&2\ten{15}\ev^{-1}&3\ten{8}\me&1\se\\
1\un{mm}&5\ten{4}\ev^{-1}&1\ten{-2}\me&3\ten{-12}\se\\
1\un{Rydberg}&2\ev&1\ten{7}\me^{-1}&3\ten{15}\se^{-1}\\
1\rb&3\ten{-4}\ev^{-1}&5\ten{-11}\me&2\ten{-19}\se\\
m_e&5\ten{6}\ev&3\ten{12}\me^{-1}&8\ten{20}\se^{-1}\\
\tau_p&3\ten{52}\ev^{-1}&5\ten{45}\me&2\ten{37}\se\\
f_\pi&1\ten{8}\ev&7\ten{14}\me^{-1}&2\ten{23}\se^{-1}\\
G_F&1\ten{-23}\ev^{-2}&5\ten{-37}\me^2&5\ten{-54}\se^2\\
G_N&7\ten{-57}\ev^{-2}&3\ten{-70}\me^2&3\ten{-87}\se^2\\
G_N^{\fr{1}{2}}G_F^{-1}&7\ten{-6}\ev&4\ten{1}\me^{-1}&1\ten{10}\se^{-1}\\
(G_F G_N)^{-\fr{1}{4}}&6\ten{19}\ev&3\ten{26}\me^{-1}&9\ten{34}\se^{-1}\\
G_N^{\fr{3}{2}}G_F^{-2}&4\ten{-39}\ev&2\ten{-32}\me^{-1}&6\ten{-24}\se^{-1}\\
r_E&3\ten{13}\ev^{-1}&6\ten{6}\me&2\ten{-2}\se\\
d_{ES}&8\ten{17}\ev^{-1}&2\ten{11}\me&5\ten{2}\se\\
d_{EM}&2\ten{15}\ev^{-1}&4\ten{8}\me&1\se\\
g&2\ten{-23}\ev&1\ten{-16}\me^{-1}&3\ten{-8}\se^{-1}\\
\tuni&7\ten{32}\ev^{-1}&1\ten{26}\me&4\ten{17}\se\\
\rhouni&4\ten{-11}\ev^4&3\ten{16}\me^{-4}&2\ten{50}\se^{-4}\\
\rhomw&8\ten{-4}\ev^{-4}&5\ten{23}\me^{-4}&4\ten{57}\se^{-4}\\
1\un{b}&3\ten{-15}\ev^{-2}&1\ten{-28}\me^2&1\ten{-45}\se^2\\
1\un{pb}&3\ten{-27}\ev^{-2}&1\ten{-40}\me^2&1\ten{-57}\se^2\\
\rhoair&6\ten{15}\ev^4&4\ten{42}\me^{-4}&3\ten{76}\se^{-4}\\
\rhowater&4\ten{18}\ev^4&3\ten{45}\me^{-4}&2\ten{79}\se^{-4}\\
\rhoau&8\ten{19}\ev^4&5\ten{46}\me^{-4}&4\ten{80}\se^{-4}}
Because $\hbar=c=1$, velocity and angular momentum are dimensionless.
\e{v_e=4\ten{-5}}
I interpret the question as explicitly asking for units in meter, second, and eV. So, for the variables containing Coulomb, the only way to convert the Coulombs is via $\epz$. With this permitted, the answers are
\e{1\un{T}=2\ten{2}\ev^2=5\ten{15}\me^{-2}=4.5\ten{32}\se^{-2}}
\e{1\un{$\mu$G}=2\ten{-8}\ev^2=5\ten{5}\me^{-2}=5\ten{22}\se^{-2}}
\e{1\ B_\textrm{Magnetar}=10^{10}\un{T}=2\ten{-12}\ev^2=5\ten{25}\me^{-2}=5\ten{42}\se^{-2}}
\e{1\un{A}=1\ten{3}\ev=6\ten{9}\me^{-1}=2\ten{18}\se^{-1}}
For the variables containing Kelvin, the only way to convert the Kelvin is via $k_B$. With this permitted, the answers are
\e{\tro=3\ten{5}\ev=1\ten{10}\me^{-1}=4\ten{18}\se^{-1}}
\e{\tcmb=2\ten{1}\ev=1\ten{8}\me^{-1}=4\ten{16}\se^{-1}}
\e{\tsun=5\ten{5}\ev=3\ten{11}\me^{-1}=8\ten{19}\se^{-1}}
Some comments.
\e{f_\pi=130\un{MeV}}
The energy density of the universe was taken from $\rho_\textrm{crit}$.
\e{\rho_{\textrm{crit}}=9\ten{-27}\un{kg}\un{m}^{-3}}
\section{Problem 2}
\begin{center}
\begin{tabular}{L L L L L L}
\textrm{Particle}&\textrm{Mass [MeV]}&\textrm{Lifetime [s]}&\textrm{Decay interaction}&\textrm{Decay process}&\textrm{Decay phase space [eV]}\\
n&939&8.79\ten{2}&\rom{Weak}&pe^{-}\bar{\nu_e}&1.04\ten{-8}\\
\Delta^{++}&1.23\ten{3}&0.56\ten{-23}&\rom{Electroweak}&\Delta^{++}\rightarrow p\pi^{+}&7.82\ten{6}\\
\Delta^{+}&1.23\ten{3}&0.56\ten{-23}&\rom{Electroweak}&\Delta^{+}\rightarrow p\pi^{0}&1.57\ten{8}\\
\Delta^{0}&1.23\ten{3}&0.56\ten{-23}&\rom{Electroweak}&\Delta^0\rightarrow n\pi^{0}&1.55\ten{8}\\
\Delta^{-}&1.23\ten{3}&0.56\ten{-23}&\rom{Electroweak}&\Delta^{-}\rightarrow n\pi^{-}&4.29\ten{8}\\
\pi^{+}&139&2.6\ten{-8}&\rom{Electroweak}&\pi^{+}\rightarrow \mu^{+}\nu_\mu&3.35\ten{7}\\
\pi^{-}&139&2.6\ten{-8}&\rom{Electroweak}&\pi^{-}\rightarrow \mu^{-}\bar{\nu_\mu}&3.35\ten{7}\\
\pi^0&134.97&8.52\ten{-17}&\rom{Electromagnetic}&\pi^{0}\rightarrow 2\gamma&1.34\ten{8}\\
\rho^0&775&4.45\ten{-23}&\rom{Strong}&\rho^0\rightarrow \pi^{+}\pi^{-}&4.97\ten{8}\\
K_L^0&497.6&1.263\ten{-7}&\rom{Electroweak}&K_L^0\rightarrow \pi^{\pm}e^{\pm}\nu_e&3.58\ten{8}\\
K_S^0&497.6&0.89\ten{-10}&\rom{Electroweak}&K_S^0\rightarrow \pi^{+}\pi^{-}&2.196\ten{8}\\
J/\psi&3096&7.08\ten{-21}&\rom{Strong}&J/\psi\rightarrow \rom{ggg}&3.09\ten{9}\\
\mu&105.5&2.19\ten{-6}&\rom{Electroweak}&\mu\rightarrow e^{-}\bar{\nu_e}\nu_\mu&1.05\ten{8}\\
\Lambda&1115&2.632\ten{-10}&\rom{Electroweak+Strong}&\Lambda\rightarrow p\pi^{-}&3.77\ten{7}
\end{tabular}
\end{center}
For this answer we will use the PDG, decay width $\Gamma=\hbar \tau^{-1}$, decay length $L=c\tau$, and 3 significant figures. For the dominant decay mode, we will use the PDG entry always. $\Delta^{++}$, $\Delta^{+}$, $\Delta^{0}$, and $\Delta^{-}$ masses are refer to the pole positions, which have real and imaginary masses, so the real part is taken. The ``dominant decay'' refers to the experimental decay mode. The dominant decay mode determines the decay mode, which is listed, and the $\Gamma$ appropriate for the decay mode is used in the calculation. The decay time $\tau$ is taken from the PDG explicitly, then from the decay rate $\Gamma$, then from the decay width $\Gamma$, whichever is first. The $J/\psi$ $\tau$ disagrees between Wikipedia and the PDG. The $K_K^0$ $\tau$ disagrees between Wikipedia and the PDG.
\section{Problem 3.a.}
\e{M=18.015\un{g}\un{mol}^{-1}}
\e{V=1\ten{6}\un{m}^3}
\e{\rho=997\un{kg}\un{m}^{-3}}
\e{N_p=\rho V M^{-1} N_a = 3.333\ten{34}\un{protons}}
\e{\tau=5.4\ten{33}\un{y}}
\e{\boxed{R=\tau^{-1}N=0.617\un{year}^{-1}}}
For the proposed detector, the expected number of events per year is 0.617.
\section{Problem 3.b.}
\e{d=36\un{m}}
\e{A=\pi r^2=1017\un{m}^2}
\e{V=A(36\un{m})=11164\pi\un{m}^3}
\e{N_p=\rho V M^{-1} N_a = 1.2213\ten{33}\un{protons}}
\e{R=\tau^{-1}N=0.225\un{year}^{-1}}
\e{\boxed{R(1000\un{day})=0.618}}
Given $\tau$, the expectation number of events inside Super-K over 1000 days is 0.618.
\section{Problem 3.c.}
If efficiency is defined as
\e{\eta=\fr{\rom{Number of observed events}}{\rom{Number of expected events}}}
then the efficiency is 0. If efficiency is defined as
\e{\eta=\fr{\rom{Number of expected events}}{\rom{Number of events at 95\% CL}}}
then the efficiency is
\e{\boxed{\eta=\fr{0.618}{3}=0.206}}
\section{Problem 3.d.}
Super-K should operate such that $N$ events are expected to appear in the detector. Because the detector has efficiency $\eta<1$, it is then expected to sample $N\eta=30$ of those events. This improves the limit on the expected number of samples by a factor of 10, from 3 to 30. Assume $\eta$ is the same as Super-K.
\e{\eta=0.206}
\e{30\eta^{-1}=145\un{events}}
The proposed detector at Super-K's efficiency is expected to observe 30 out of 145 total events in the detector. The expected number of years for 145 events to appear is
\e{\boxed{(30 \eta^{-1})(\tau^{-1} N)^{-1}=236\un{y}}}
\section{Problem 4}
The question is ambiguous as to what the $U(1)$ transformations are. Assume they are global transformations,
\e{U(1)_A:\quad \psi\rightarrow e^{i\alpha \gamma^5} \psi}
\e{U(1)_V:\quad \psi\rightarrow e^{i\alpha} \psi}
\ctab{L L L L L}{
\rom{Term} & \rom{$\ua$ invariant} & \rom{$\uv$ invariant} & \rom{Lorentz invariant} & \rom{Renormalizable}\\
\psb i\sib^m \p_m \psi & \yes & \no & \yes&\yes\\
\p_m \psb \p^m \psi &\yes&\no&\yes&\yes\\
\psi^c \p_m \p^m \psi&\no&\no&\no&\yes\\
\psb \psi&\yes&\no&\yes&\yes\\
\psi \psi^c&\no&\yes&\no&\yes\\
\psi \psi&\no&\no&\no&\yes\\
\psb \sigma^{mn}\psi F_{mn}&\yes&\no&\yes&\yes\\
\psi \sigma^{mn} \psi F_{mn}&\no&\no&\no&\yes\\
\psi^c \sigma^{mn} \psi F_{mn}&\no&\yes&\no&\yes\\
\psb \sib^m \psi \psb \sib_m \psi&\yes&\no&\yes&\no\\
\psi^c \psi \psi^c \psi&\no&\yes&\no&\no\\
\psi^c \sigma^{mn} \psi \psb^c \sib_{mn} \psb&\yes&\no&\no&\no
}
By applying the transformation rules for $\ua$ and $\uv$ on $\psi$, $\psb$, $\psi^c$, and $\bar{\psi}^c$ we determined if the terms are invariant and renormalizable.
\end{document}
